El estado del arte se ha elaborado tomando como referencia las herramientas que ya forman parte del entorno del distribuidor y, además, soluciones comerciales existentes que abordan alguna de las necesidades detectadas en la Sección~\ref{sec:context}. El objetivo no es localizar una aplicación idéntica al sistema propuesto, sino comparar alternativas reales que cubren parcialmente ámbitos como la gestión de clientes, la venta B2B, la preventa en ruta, el software de salón, las aplicaciones web progresivas y la planificación de rutas.

\subsection{Tecnologías conocidas por el distribuidor}

Durante el proceso de \textit{captación de requisitos} descrito en la Sección~\ref{sec:RequirementsCaptureMethodology}, se ha identificado que el distribuidor tiene experiencia previa con herramientas de gestión empresarial y administrativa. Las tecnologías conocidas pueden agruparse de la siguiente forma:

\begin{enumerate}[label=\textbf{\arabic*.}, leftmargin=1.1cm, itemsep=0.45em, topsep=0.25em]
  \item \textbf{Herramientas ofimáticas y bases de datos locales:}
  \begin{itemize}[topsep=0.1em, itemsep=0.15em]
    \item \textbf{Microsoft Excel:} herramienta utilizada para la \textbf{gestión de pedidos}, el seguimiento de inventarios, la elaboración de informes comerciales y el mantenimiento de listados operativos.
    \item \textbf{Microsoft Access:} solución conocida por el distribuidor para estructurar información en bases de datos locales y apoyar procesos internos como expediciones o listados de reparto.
    \item \textbf{Google Sheets:} hoja de cálculo en línea utilizada como alternativa flexible para registrar información operativa y compartir datos de manera sencilla cuando no se requiere una aplicación específica.
  \end{itemize}

  \item \textbf{Software de facturación y gestión administrativa:}
  \begin{itemize}[topsep=0.1em, itemsep=0.15em]
    \item \textbf{TeamSystem Factusol:} software empleado para cubrir procesos propios de \textbf{facturación}, cobros y documentación comercial dentro del contexto administrativo del distribuidor.
    \item \textbf{TeamSystem Preventa:} solución del mismo ecosistema orientada a la \textit{autoventa} y \textit{preventa} en movilidad. Permite crear presupuestos, pedidos, albaranes y facturas en ruta, gestionar cobros y transferir información hacia TeamSystem Delsol o TeamSystem Factusol \citep{sdelsol-preventa}.
  \end{itemize}
\end{enumerate}

Estas herramientas resultan útiles porque ya son conocidas por el negocio y resuelven tareas concretas. Sin embargo, desde el punto de vista de los \textbf{sistemas de información} presentan limitaciones importantes: la información queda repartida entre varias fuentes, no existe una experiencia única para clientes, comerciales y administradores, y la trazabilidad depende en gran medida de procedimientos manuales o de sincronizaciones externas. Por tanto, aunque el sistema propuesto debe poder convivir con este entorno, su finalidad no es reproducir una hoja de cálculo ni sustituir por completo a un \textit{ERP}, sino centralizar la operativa comercial específica del distribuidor.

\subsection{Aplicaciones existentes de la empresa central}

La empresa central dispone de recursos digitales propios que resultan relevantes para el dominio del proyecto, especialmente por su relación con el catálogo de productos y la formación técnica. Las aplicaciones y canales identificados son los siguientes:

\begin{enumerate}[label=\textbf{\arabic*.}, leftmargin=1.1cm, itemsep=0.45em, topsep=0.25em]
  \item \textbf{\textit{KIN Cosmetics} como aplicación móvil oficial:}
  \begin{itemize}[topsep=0.1em, itemsep=0.15em]
    \item Se encuentra disponible para dispositivos móviles y actúa principalmente como \textbf{catálogo digital de productos} y como repositorio de contenidos formativos para profesionales del sector \citep{kintrends}.
    \item Permite consultar las distintas líneas de producto, acceder a información técnica y visualizar material formativo relacionado con técnicas de peluquería.
    \item Su aportación principal se sitúa en la \textbf{consulta de producto} y en la \textbf{formación técnica}, no en la gestión diaria de la red comercial.
    \item No ofrece funcionalidades orientadas a la gestión operativa del distribuidor, tales como la planificación de visitas comerciales, la gestión de pedidos propios, el seguimiento de entregas, la segmentación de clientes o la integración con procesos administrativos.
  \end{itemize}

  \item \textbf{Página web oficial de la empresa:}
  \begin{itemize}[topsep=0.1em, itemsep=0.15em]
    \item Permite consultar el catálogo de productos e incluye información detallada sobre sus características \citep{kincosmetics_web}.
    \item En algunos casos contempla la posibilidad de realizar pedidos, aunque el proceso no se encuentra integrado dentro de un sistema propio del distribuidor.
    \item La compra se basa en la redirección a plataformas externas de comercio electrónico, por lo que no queda incorporada a una operativa propia de clientes, comerciales y administración local.
  \end{itemize}

  \item \textbf{Plataformas externas de comercio electrónico:}
  \begin{itemize}[topsep=0.1em, itemsep=0.15em]
    \item El producto \textit{KINESSENCES Nourish Intense Mask}, por ejemplo, puede visualizarse en la web oficial \citep{kincosmetics_producto}, desde donde se facilita el acceso a su compra a través de plataformas como Amazon \citep{amazon_producto}.
    \item Este enfoque facilita la adquisición puntual de productos, pero no proporciona herramientas específicas para gestionar la relación continuada entre distribuidor, comerciales y peluquerías clientes.
    \item Tampoco permite modelar el \textbf{ciclo operativo completo del negocio}: visita, asesoramiento, pedido, reparto, cobro, comunicación y seguimiento técnico del salón.
  \end{itemize}
\end{enumerate}

\subsection{Soluciones alternativas existentes en el mercado}

Además de las herramientas ya conocidas por el distribuidor, existen soluciones comerciales consolidadas que cubren partes del problema. A continuación se analizan por familias funcionales, atendiendo a su utilidad potencial y a sus limitaciones respecto al contexto concreto del proyecto.

\begin{enumerate}[label=\textbf{\arabic*.}, leftmargin=1.1cm, itemsep=0.6em, topsep=0.35em]
  \item \textbf{CRM generalistas:}
  \begin{itemize}[topsep=0.1em, itemsep=0.15em]
    \item \textbf{HubSpot CRM:} plataforma orientada a centralizar datos de clientes, contactos, operaciones, tareas y paneles de informes \citep{hubspot-crm}. Su utilidad principal se sitúa en la gestión de relaciones comerciales y oportunidades de venta.
    \item \textbf{Zoho CRM:} solución que destaca la automatización de actividades a lo largo del embudo comercial, la definición de procesos mediante flujos guiados y la personalización de módulos y diseños \citep{zoho-crm-features}.
  \end{itemize}

  Estas soluciones aportan una base sólida para equipos de ventas de propósito general, pero no están diseñadas específicamente para la \textbf{distribución de productos capilares a peluquerías}. Requieren configuración o desarrollo adicional para cubrir catálogo técnico, pedidos B2B recurrentes, rutas de reparto, control operativo de entregas, fichas técnicas del salón y sincronización con herramientas administrativas ya usadas por el distribuidor.

  \item \textbf{ERP ligero, facturación y preventa:}
  \begin{itemize}[topsep=0.1em, itemsep=0.15em]
    \item \textbf{Odoo Sales:} solución de gestión empresarial que permite trabajar con presupuestos, pedidos de venta, facturación, portal de cliente, tarifas, descuentos, productos e integración con módulos como inventario, contabilidad y comercio electrónico \citep{odoo-sales-features}.
    \item \textbf{TeamSystem Delsol y Factusol:} ecosistema más próximo a la dimensión administrativa del problema, especialmente por su relación con facturación, cobros y documentación comercial.
    \item \textbf{TeamSystem Preventa:} herramienta especialmente relevante por estar orientada a comerciales en ruta. Contempla documentos de venta, catálogo de imágenes, consulta de stock, cobros, visitas y conexión con Factusol \citep{sdelsol-preventa}.
  \end{itemize}

  Estas herramientas cubren procesos empresariales de forma más completa que una hoja de cálculo, pero su adopción implica adaptar el negocio a un producto existente y, en algunos casos, configurar flujos que exceden el alcance real del distribuidor. Además, el proyecto desarrollado no pretende convertirse en un \textit{ERP} generalista ni sustituir la facturación oficial, sino proporcionar una capa operativa propia para usuarios internos y clientes profesionales, dejando la integración completa con Factusol como una evolución posterior.

  \item \textbf{Plataformas de pedidos B2B:}
  \begin{itemize}[topsep=0.1em, itemsep=0.15em]
    \item \textbf{Shopify B2B:} plataforma que permite definir catálogos y precios distintos según empresa o ubicación, limitar productos disponibles y aplicar reglas de cantidad o precios por volumen \citep{shopify-b2b-catalogs}.
    \item \textbf{Pepperi:} solución orientada específicamente a marcas y distribuidores, combinando comercio electrónico B2B, toma de pedidos móvil para representantes, promociones, ejecución comercial, analítica operativa y capacidades de \textit{Direct Store Delivery} \citep{pepperi-platform}.
  \end{itemize}

  Este tipo de soluciones se acerca mucho a la problemática del \textbf{pedido profesional} y de la \textbf{venta multicanal}. Sus puntos fuertes son el catálogo, la personalización de precios y la toma de pedidos por comerciales o compradores. Sin embargo, suelen estar pensadas para organizaciones con procesos de venta B2B más estandarizados y no resuelven por sí solas la gestión técnica del salón, las necesidades concretas de formación y comunicación del distribuidor, ni la integración gradual con el sistema administrativo ya existente. Además, pueden introducir costes de licencia, dependencia de proveedor y complejidad de implantación superiores a los asumibles en el contexto del TFG.

  \item \textbf{Software de salón y reservas:}
  \begin{itemize}[topsep=0.1em, itemsep=0.15em]
    \item \textbf{Booksy:} aplicación orientada a la gestión de salones que incorpora funciones de calendario, cobros, marketing por texto y correo, etiquetas de clientes, gestión de personal, informes e inventario en algunos planes \citep{booksy}.
    \item \textbf{Fresha:} solución para negocios de belleza y bienestar que ofrece planificación de citas, pagos, gestión de clientes, recordatorios automatizados, formularios, punto de venta y funcionalidades de marketing \citep{fresha-features}.
  \end{itemize}

  Estas aplicaciones son relevantes porque se dirigen al mismo entorno profesional de peluquerías y salones. Sin embargo, su usuario principal es el propio salón, no el distribuidor de productos. Su núcleo funcional se centra en reservas, agenda, caja y relación entre salón y cliente final. Por ello, no cubren de forma natural la gestión de rutas comerciales del distribuidor, el pedido de producto profesional entre empresas, el reparto físico, la relación comercial con múltiples peluquerías o la coordinación con sistemas de facturación externos.

  \item \textbf{Aplicaciones web progresivas comerciales:}
  \begin{itemize}[topsep=0.1em, itemsep=0.15em]
    \item \textbf{\textit{web.dev} y el enfoque PWA:} las aplicaciones web progresivas no constituyen por sí mismas una solución de negocio, sino una forma de entregar aplicaciones web con capacidades cercanas a las aplicaciones nativas. Según la documentación de \textit{web.dev}, una PWA puede ser instalable, ejecutarse en una ventana independiente, funcionar sobre una única base de código y ofrecer ciertas capacidades sin conexión cuando el navegador y la plataforma lo permiten \citep{webdev-pwa}.
    \item \textbf{Adobe Commerce PWA Studio:} herramienta orientada a desarrollar escaparates PWA sobre Adobe Commerce o Magento Open Source dentro del ámbito del comercio electrónico \citep{adobe-pwa-studio}.
  \end{itemize}

  Este enfoque resulta adecuado porque los usuarios necesitan acceder desde móvil durante visitas comerciales, reparto o consulta en el salón, sin imponer la instalación desde tiendas de aplicaciones. Sin embargo, una PWA comercial o un escaparate de comercio electrónico no aporta automáticamente el modelo de datos, los roles, la trazabilidad, la planificación de rutas o los flujos operativos específicos. La \textit{PWA} es, por tanto, una decisión tecnológica de entrega, pero debe combinarse con un diseño funcional propio.

  \item \textbf{Sistemas de gestión de rutas y última milla:}
  \begin{itemize}[topsep=0.1em, itemsep=0.15em]
    \item \textbf{Onfleet:} herramienta especializada en optimización de rutas que tiene en cuenta ubicaciones de entrega, ventanas horarias, horarios de conductores, capacidades de vehículos, tiempos de servicio, tráfico histórico y restricciones de carretera \citep{onfleet-route-optimization}.
    \item \textbf{OptimoRoute:} solución orientada a planificar y optimizar rutas para múltiples conductores y paradas, considerando restricciones como ventanas temporales, horarios, capacidades y límites de vehículos \citep{optimoroute-features}.
    \item \textbf{Routific y Route4Me:} herramientas identificadas dentro del mismo grupo de soluciones especializadas en planificación de rutas y entregas de última milla.
  \end{itemize}

  Estas soluciones son potentes para empresas cuya dificultad principal es la \textbf{optimización logística}. En el caso del distribuidor analizado, la ruta es una parte importante del proceso, pero no el único núcleo del sistema. La aplicación desarrollada necesita vincular rutas con clientes, visitas, pedidos, entregas, comunicación, catálogo y administración. Una herramienta de última milla podría complementar el sistema en una fase avanzada, pero no sustituye a una plataforma integrada de gestión comercial.
\end{enumerate}

\subsection{Resumen y conclusiones de la comparativa}

La Tabla~\ref{tab:comparativa-alternativas-mercado} resume los aspectos positivos y las principales limitaciones de las familias de soluciones analizadas respecto a las necesidades del distribuidor.

{\small
\renewcommand{\arraystretch}{1.22}
\begin{table}[H]
\centering
\caption{Comparativa de alternativas existentes en el mercado}
\label{tab:comparativa-alternativas-mercado}
\begin{tabularx}{\textwidth}{>{\RaggedRight\arraybackslash}p{0.20\textwidth}>{\RaggedRight\arraybackslash}p{0.34\textwidth}>{\RaggedRight\arraybackslash}p{0.34\textwidth}}
\hline
\textbf{Alternativa} & \textbf{Aspectos positivos} & \textbf{Limitaciones respecto al proyecto} \\
\hline
Hojas de cálculo y bases de datos locales & Flexibles, conocidas por el distribuidor, bajo coste inicial y útiles para listados, cálculos y registros sencillos. & Información fragmentada, duplicidad de datos, ausencia de roles, baja trazabilidad y dificultad para ofrecer acceso estructurado a comerciales y clientes. \\
\hline
KIN Cosmetics y web oficial de la marca & Catálogo oficial, contenido técnico, material formativo y referencia directa de producto. & Orientación corporativa y de consulta; no cubre pedidos propios del distribuidor, rutas, cobros, visitas, comunicación segmentada ni administración local. \\
\hline
CRM generalistas: HubSpot, Zoho & Buena gestión de contactos, oportunidades, tareas, automatización comercial, informes y seguimiento de relaciones. & No modelan de serie catálogo capilar técnico, pedidos B2B con entrega local, fichas técnicas de salón ni integración específica con Factusol. \\
\hline
ERP y preventa: Odoo, Factusol, TeamSystem Preventa & Cubren presupuestos, pedidos, facturas, inventario, cobros, tarifas y trabajo comercial en movilidad. & Mayor complejidad de implantación; el foco está en gestión empresarial general o documentación comercial, no en una PWA específica para distribuidor, comerciales y peluquerías. \\
\hline
\end{tabularx}
\end{table}

\begin{table}[H]
\centering
\ContinuedFloat
\caption[]{Comparativa de alternativas existentes en el mercado (continuación)}
\begin{tabularx}{\textwidth}{>{\RaggedRight\arraybackslash}p{0.20\textwidth}>{\RaggedRight\arraybackslash}p{0.34\textwidth}>{\RaggedRight\arraybackslash}p{0.34\textwidth}}
\hline
\textbf{Alternativa} & \textbf{Aspectos positivos} & \textbf{Limitaciones respecto al proyecto} \\
\hline
Pedidos B2B: Shopify B2B, Pepperi & Digitalizan catálogos, precios por cliente, pedidos profesionales, promociones y venta multicanal. & Pueden resultar sobredimensionadas o dependientes de proveedor; no resuelven por sí solas la operativa técnica del salón, rutas locales y adaptación al proceso actual. \\
\hline
Software de salón: Booksy, Fresha & Muy adecuados para agenda, reservas, clientes del salón, recordatorios, pagos, marketing y punto de venta. & Su usuario objetivo es el salón frente a su cliente final, no el distribuidor frente a peluquerías; no cubren reparto, preventa, pedidos a proveedor ni control administrativo del distribuidor. \\
\hline
PWA comerciales y escaparates PWA & Permiten acceso móvil, instalación desde navegador, una única base de código y experiencias cercanas a una app nativa. & Son una forma de entrega tecnológica, no una solución funcional completa; requieren diseñar el dominio, roles, datos y procesos propios. \\
\hline
Rutas y última milla: Onfleet, OptimoRoute, Routific, Route4Me & Optimizan rutas con ventanas horarias, capacidades, seguimiento, despacho y restricciones logísticas. & Cubren la capa logística, pero no integran catálogo, pedidos, clientes profesionales, comunicación, fichas técnicas ni administración del distribuidor. \\
\hline
\end{tabularx}
\end{table}
}

El análisis anterior muestra que existen alternativas reales y maduras, pero su cobertura es parcial respecto al problema abordado. Las herramientas de oficina resuelven registros simples, las aplicaciones de la marca central ofrecen catálogo y formación, los \textit{CRM} generalistas organizan la relación comercial, los \textit{ERP} y soluciones de preventa cubren documentación administrativa, las plataformas B2B facilitan pedidos entre empresas, el software de salón gestiona reservas y clientes finales, y los sistemas de rutas optimizan la última milla.

La \textbf{principal limitación} es que ninguna de estas alternativas integra, dentro de una única experiencia adaptada al distribuidor, los procesos de gestión de usuarios, clientes profesionales, catálogo técnico, pedidos, visitas comerciales, reparto, notificaciones, promociones, formaciones, trazabilidad administrativa y soporte a PWA. Adoptar una herramienta generalista exigiría adaptar el proceso del distribuidor al producto elegido o desarrollar integraciones adicionales, mientras que usar varias herramientas a la vez mantendría la fragmentación de información descrita en la Sección~\ref{sec:context}.

Por este motivo, el sistema propuesto se justifica como una \textbf{solución específica y acotada} al dominio del distribuidor. Su valor no reside en competir con un \textit{ERP} completo, un \textit{CRM} comercial o una plataforma logística especializada, sino en integrar las funciones necesarias para el caso real en una aplicación web progresiva con roles diferenciados, modelo de datos centralizado y trazabilidad suficiente para evolucionar posteriormente hacia integraciones más avanzadas.
